\documentclass[a4paper]{article}

\newcommand{\docTitle}{Verteilte Systeme}

\usepackage{datetime2}  % for dates
\usepackage[utf8]{inputenc} % UTF-8 input encoding
\usepackage[T1]{fontenc} % fixes \textquotedbl with pdflatex/OT1
\usepackage{textcomp}   % additional text symbols for typewriter/upquote
\usepackage{xcolor}     % for gray color
\usepackage{listings}
\usepackage{marginnote} % for horizontal margin notes
\usepackage{parskip}
\usepackage{amsmath,amssymb,mathtools}
\usepackage{everypage}
\usepackage{fancyhdr}
\usepackage{amsfonts}   % Mengen and natural numbers
\usepackage{enumerate}  % Custom Enumerations
\usepackage[inline]{enumitem}   % List spacing and labels
\usepackage{multicol}   % Multiple Columns
\usepackage{tabularx}

\usepackage{hyperref}   % Links
\hypersetup{
    pdftitle={\docTitle},
    pdfauthor={Moritz Köhler},
    pdfkeywords={DHBW, Hochschule},
    colorlinks=true,
    linkcolor=black,
    filecolor=magenta,
    urlcolor=cyan
}

% -------
% Images
% -------

\usepackage{import}
\usepackage{xifthen}
\usepackage{pdfpages}
\usepackage{tikz}
\usetikzlibrary{arrows.meta, positioning}

\newcommand{\incfig}[1]{%
    \def\svgwidth{\columnwidth}
    \import{./fig/}{#1.pdf_tex}
}

% ---------------------------------
% Vertical Side note (Bottom Left)
% ---------------------------------

\newcommand{\verticaldocinfo}{%
  \rotatebox{90}{\tiny\color{gray}\texttt{\DTMtoday\_\docTitle\_\jobname}}%
}

\AddEverypageHook{
  \put(-40pt,-650pt){\verticaldocinfo}%
}

% ---------------
% Math Shortcuts
% ---------------

% Number sets
\newcommand{\R}{\mathbb{R}}
\newcommand{\N}{\mathbb{N}}
\newcommand{\Z}{\mathbb{Z}}
\newcommand{\Q}{\mathbb{Q}}
\newcommand{\C}{\mathbb{C}}

% Vectors and matrices
\newcommand{\vect}[1]{\mathbf{#1}}
\newcommand{\mat}[1]{\mathbf{#1}}

% Derivatives
\newcommand{\dd}{\mathrm{d}}
\newcommand{\dpartial}[2]{\frac{\partial #1}{\partial #2}}
\newcommand{\ddt}{\frac{\dd}{\dd t}}

% Norms and absolute values
\newcommand{\norm}[1]{\left\lVert #1 \right\rVert}
\newcommand{\abs}[1]{\left\lvert #1 \right\rvert}

% Probability & Statistics
\newcommand{\E}{\mathbb{E}}
\newcommand{\Var}{\mathrm{Var}}
\newcommand{\Prob}{\mathbb{P}}

% Fields and Algebras
\newcommand{\K}{\mathbb{K}}
\newcommand{\F}{\mathbb{F}}

% Set Operations
\newcommand{\compl}[1]{#1^{\mathrm{c}}}

% Matrix and Vector Operations
\newcommand{\T}{^{\top}}
\newcommand{\inv}{^{-1}}
\newcommand{\conj}[1]{\overline{#1}}

% -------------------
% Main lecture macro 
% -------------------

\newcommand{\lecture}[1]{%
  \protect\marginnote{\protect\tiny\protect\color{gray}#1}
}

% -----------------------------
% Command for box with heading
% -----------------------------

\fboxsep2mm
\newcommand{\know}[2]{
    \noindent\fbox{
        \parbox{\textwidth-21pt}{

            \textbf{#1:}
            \vspace{0.125cm}

            #2
        }
    }
}

% -------------------
% Listings Languages
% -------------------

% Shared defaults for code listings
\lstset{
    basicstyle=\ttfamily\small,
    keywordstyle=\color{blue}\bfseries,
    commentstyle=\color{gray},
    stringstyle=\color{teal},
    numberstyle=\tiny\color{gray},
    numbers=left,
    stepnumber=1,
    numbersep=8pt,
    showstringspaces=false,
    frame=single,
    breaklines=true,
    tabsize=4,
    keepspaces=false,
    columns=flexible,
    upquote=true,
    extendedchars=false % disable extended char handling to avoid UTF-8 conflicts
}

% SQL syntax highlighting
\lstdefinelanguage{SQL}{
    morekeywords={SELECT,FROM,WHERE,INSERT,INTO,VALUES,UPDATE,SET,DELETE,CREATE,TABLE,PRIMARY,KEY,NOT,NULL,AND,OR,JOIN,LEFT,RIGHT,INNER,ON,AS,ORDER,BY,GROUP,HAVING,LIMIT},
    sensitive=false,
    morecomment=[l]{--},
    morestring=[b]'
}

% JSON syntax highlighting
\lstdefinelanguage{json}{
    morestring=[b]",
    stringstyle=\color{teal},
    comment=[l]{//},
    morecomment=[s]{/*}{*/},
        morekeywords={true,false,null},
    literate=
     *{0}{{{\color{blue}0}}}{1}
      {1}{{{\color{blue}1}}}{1}
      {2}{{{\color{blue}2}}}{1}
      {3}{{{\color{blue}3}}}{1}
      {4}{{{\color{blue}4}}}{1}
      {5}{{{\color{blue}5}}}{1}
      {6}{{{\color{blue}6}}}{1}
      {7}{{{\color{blue}7}}}{1}
      {8}{{{\color{blue}8}}}{1}
      {9}{{{\color{blue}9}}}{1}
            {:}{{{\color{black}:}}}{1}
            {,}{{{\color{black},}}}{1},
    showstringspaces=false
}

% Language specific styles
\lstdefinestyle{javastyle}{
    language=Java,
    basicstyle=\ttfamily\small,
    keywordstyle=\color{blue}\bfseries,
    commentstyle=\color{gray},
    stringstyle=\color{teal},
    numberstyle=\tiny\color{gray},
    numbers=left,
    stepnumber=1,
    numbersep=8pt,
    showstringspaces=false,
    frame=single,
    breaklines=true,
    tabsize=4,
    keepspaces=true,
    columns=fullflexible,
    upquote=true
}

\lstdefinestyle{sqlstyle}{
    language=SQL,
    basicstyle=\ttfamily\small,
    keywordstyle=\color{blue}\bfseries,
    commentstyle=\color{gray},
    stringstyle=\color{teal},
    numberstyle=\tiny\color{gray},
    numbers=left,
    stepnumber=1,
    numbersep=8pt,
    showstringspaces=false,
    frame=single,
    breaklines=true,
    tabsize=4,
    keepspaces=true,
    columns=fullflexible,
    upquote=true
}

\lstdefinestyle{pythonstyle}{
    language=Python,
    basicstyle=\ttfamily\small,
    keywordstyle=\color{blue}\bfseries,
    commentstyle=\color{gray},
    stringstyle=\color{teal},
    numberstyle=\tiny\color{gray},
    numbers=left,
    stepnumber=1,
    numbersep=8pt,
    showstringspaces=false,
    frame=single,
    breaklines=true,
    tabsize=4,
    keepspaces=true,
    columns=fullflexible,
    upquote=true
}

\lstdefinestyle{cstyle}{
    language=C,
    basicstyle=\ttfamily\small,
    keywordstyle=\color{blue}\bfseries,
    commentstyle=\color{gray},
    stringstyle=\color{teal},
    numberstyle=\tiny\color{gray},
    numbers=left,
    stepnumber=1,
    numbersep=8pt,
    showstringspaces=false,
    frame=single,
    breaklines=true,
    tabsize=4,
    keepspaces=true,
    columns=fullflexible,
    upquote=true
}

\lstdefinestyle{bashstyle}{
    language=bash,
    basicstyle=\ttfamily\small,
    keywordstyle=\color{blue}\bfseries,
    commentstyle=\color{gray},
    stringstyle=\color{teal},
    numberstyle=\tiny\color{gray},
    numbers=left,
    stepnumber=1,
    numbersep=8pt,
    showstringspaces=false,
    frame=single,
    breaklines=true,
    tabsize=4,
    keepspaces=true,
    columns=fullflexible,
    upquote=true
}

\lstdefinestyle{jsonstyle}{
    language=json,
    basicstyle=\ttfamily\small,
    keywordstyle=\color{blue}\bfseries,
    commentstyle=\color{gray},
    stringstyle=\color{teal},
    numberstyle=\tiny\color{gray},
    numbers=left,
    stepnumber=1,
    numbersep=8pt,
    showstringspaces=false,
    frame=single,
    breaklines=true,
    tabsize=4,
    keepspaces=true,
    columns=fullflexible,
    upquote=true
}

% Default listing style
\lstset{language={}}

\title{\docTitle}
\author{Moritz}
\date{\today}

\begin{document}
\maketitle
\tableofcontents

\lecture{2026-10-07}

\section{Einführung}

\subsection{Agenda}

\begin{itemize}
    \item The World of Distributed Systems
    \item Architecture of Distributed Systems
    \item Communication (RMI, RPC)
    \item Global Sates and Synchronization
    \item (Consistency and Replication)
    \item (Fault Tolerance)
    \item Components and "their World"
\end{itemize}

Wir haben zwei Räume. So 50/50 für Theorie und Praxis.

Theorie wird meist früher enden und bei der Praxis/ Übungen machen wir auch mal bis zu einem Abschluss auch mal länger.

\subsection{Prüfungsleistung}

60\% Eine 45min Klausur ohne Unterlagen am 15.12. Probeklausur gibt es nicht und Klausuraufgaben gibt es nicht.

40\% Eine Hausarbeit in zweier Gruppen. ca. 7 S. Inhalt am 17.12. Thema dürfen wir selber auswählen. Darf auch bis Ende der Klausurwoche abgegeben wird.

Diese Prüfungsleistung wird 50/50 mit IT Architektur verrechnet und nur gemeinsam bestanden.

Mögliche Themen:

\begin{itemize}
    \item Thor Browser
    \item 
\end{itemize}

Sollen eine Liste Pflegen, damit keine Dopplungen existieren.

\section{The World of Distributed Systems}

Ein nicht verteiltes System ist ein ein-Prozess System, welches auch mehrere Cores benutzen darf.

Mittlerweile sind viele Systeme aber schon Verteilte Systems, da bei einem CLoud Sync ja auch schon zwischen Server notwendig sind.

\subsection{History}

Die Zeitgeschichte (Liste S.10) wird nicht abgefragt.

\subsection{Characterization of Distributed Systems}

Es gibt mehrere Herausforderungen:

\begin{itemize}
    \item Heterogeneity (Unterschiedliche Hardware, Schnittstellen)
    \item Openness (Ist es Erweiterbar?)
    \item Security (Confidentiality, Integrity, Availability)
    \item Scalability (Vertical and horizontal)
    \begin{itemize}
        \item Controlling the cost
        \item Controlling performance loss
        \item Avoiding bottlenecks
    \end{itemize}
    \item Failure Handling 
    \begin{itemize}
        \item Detecting failures
        \item Masking failures
        \item Tolerating failures
        \item Recovery of failures
        \item Redundancy
    \end{itemize}
    \item Concurrency (Wie wird mit gleichzeitigem Zugriff umgegangen?)
    \item Transparency
    \begin{itemize}
        \item Access Transparency
        \item Location Transparency
        \item Concurrency Transparency
        \item Replication Transparency
        \item Failure Transparency
        \item Mobility Transparency
        \item performance Transparency
        \item Scaling Transparency
    \end{itemize}
\end{itemize}

\subsection{Concepts}

\subsubsection{Hardware}

Miltiprocessor, Multicomputer, Bus-base (Cable TV), Switch-based (Telephone Network)

Homogeneous Multicomputer Systems and Heterogenous Multicomputer Systems

Switched vs. BUS

Grids vs. Hypercubes (Optimizations for their respecting problems)

\subsubsection{Software}

Loosely coupled systems

Tight coupled systems

\know{Typische Klausurfrage}{Was ist der Hauptunterschied zwischen NOS und DOS?}

\begin{tabular}{ll}
    System & Main Goal\\
    \hline
    DOS & Manage and hide hardware resources\\
    NOS & Provide local services for remote clients
\end{tabular}

\subsubsection{Network}

\href{http://en.wikipedia.org/wiki/MOSIX}{MOSIX}

% Hab nicht alles mitgeschrieben. Gilt für alle Konzepte.

\end{document}